\section{Praktik Proyek: Aplikasi Interaktif}

Menggabungkan materi bab ini dalam \textbf{proyek praktik} adalah cara terbaik untuk memperdalam pemahaman. Proyek harus melibatkan: manipulasi DOM, event handling, validasi form, komunikasi API (Fetch), dan (opsional) penggunaan library. Berikut beberapa ide proyek dengan tingkat kesulitan bertahap \cite{mdnJsGuide}.

\begin{table}[htbp]
\centering
\begin{tabular}{|P{3.5cm}|P{3.5cm}|P{5cm}|}
\hline
\textbf{Proyek} & \textbf{Konsep Utama} & \textbf{Fitur} \\
\hline
Todo List & DOM, event, array & CRUD, filter, local storage \\
\hline
Weather App & Fetch, async/await, API & Input kota, tampil cuaca \\
\hline
Quiz App & Array objek, timer, DOM & Pilihan ganda, skor, countdown \\
\hline
Form Registrasi & Validasi, regex, event & Real-time, multi-field \\
\hline
Dashboard & Fetch, charting, layout & Data dinamis, responsif \\
\hline
\end{tabular}
\caption{Ide proyek praktik JavaScript dengan konsep yang diterapkan}
\end{table}

\begin{konsep}
\textbf{Metodologi Proyek.} Mulai dengan perencanaan: (1) daftar fitur (user stories); (2) desain UI (wireframe sederhana); (3) struktur data (array/objek apa yang dibutuhkan); (4) implementasi bertahap (satu fitur per iterasi); (5) testing dan debugging; (6) refactoring. Gunakan Git untuk version control dari awal. Kode yang ``bekerja'' bukan akhir---kode yang bersih, terorganisir, dan teruji adalah tujuan \cite{mdnToolsTesting}.
\end{konsep}

